
\section{模型假设}

为简化问题，本文做出以下假设。

\begin{enumerate}[label=\textbf{假设\arabic*}, itemindent=2em, leftmargin=4em]
\item 蔬菜当日未售出则隔日无法再售，每日为一独立库存周期。该假设为题目原话，是报童模型单周期框架的基础。

\item 定价采用成本加成定价方法$P=(1+\alpha)C$。该假设依据题目明确要求，加成率$\alpha$为决策变量。

\item 打折为被动清仓行为，折扣回收率$k=0.66$。$k$从打折交易数据中估算（折扣价与正价比值的中位数），论文同时报告$k\in[0.50,0.75]$的对应区间。

\item 历史销售数据可代表未来需求趋势。该假设为所有时间序列预测模型的基础，通过留出集验证和跨年稳定性分析进行检验。

\item 损耗率在预测期内恒定。题目仅提供近期盘点损耗率，模型将其视为常数。

\item 退货可忽略。退货仅461条（占总记录的0.05\%），数据清洗时已剔除。

\item Q2以品类为单位不考虑销售空间约束。题目仅在问题三提出销售空间限制。
\end{enumerate}
